% GENERATED FILE. Edit canonical lessons, then run npm run generate:books.
% canonical-source-hash: fnv1a64:a9cc5e4240e435ed
% canonical-lessons: GU-C512-pad, GU-C512-intaravyu, GU-C512-vetan, GU-C512-karkirdi, GU-C512-badhti

\chapter{Looking for a Job}
\label{ch:gu-looking-for-a-job}
\hlchaptermodality{512}

\begin{chapteropening}
\textbf{By the end of this chapter:} \emph{I can talk about candidates, interviews, pay and promotion.}

Complete the last lesson of chapter 512: I can talk about candidates, interviews, pay and promotion.
\end{chapteropening}

\section[\texorpdfstring{\gu{પદ} (pad)}{પદ (pad)}]{\gu{પદ} (pad) — a job post}

\label{lesson:GU-C512-pad}



\begin{quote}
Before the new one: say the Gujarati for a mechanic, then the Gujarati for a clerk.
\end{quote}

\subsection*{You'll want to know: \gu{પદ}}
\textbf{\gu{પદ}} — \emph{pad} — \textquotedblleft{}a job post\textquotedblright{}.

\begin{quote}
\gu{આ} \gu{પદ} \gu{માટે} \gu{ઘણી} \gu{અરજીઓ} \gu{આવી}.
\end{quote}

\begin{quote}
\emph{ā pad māṭe ghaṇī arjīo āvī.}
\end{quote}

\begin{quote}
Many applications came in for this post.
\end{quote}

A neuter noun.

\begin{cousinweb}
\textbf{\gu{પદ}} (\emph{pad}) is a tatsama, taken from Sanskrit \emph{pada}, \textquotedblleft{}a step, a footing; a position\textquotedblright{} (Monier-Williams); in an office it is a post one is appointed to.
\end{cousinweb}

\begin{grammarlens}[title={What you've built}]
One more word for describing what you lived and felt.
\end{grammarlens}

\subsection*{Your turn}
\begin{itemize}
\raggedright
  \item \emph{Say it:} \emph{pad}
  \item \emph{Say it:} \emph{pad}, once more
  \item \emph{Say it:} \emph{ā pad māṭe ghaṇī arjīo āvī.} — Many applications came in for this post.
  \item \emph{Recall:} say the Gujarati for a mechanic, then the Gujarati for a clerk, then say \emph{pad} again
\end{itemize}

\begin{tcolorbox}[breakable,colback=teal!4,colframe=teal!35!black,title={Before you move on}]
What does \gu{પદ} mean? (\textquotedblleft{}A job post\textquotedblright{}.) Say it once more.
\end{tcolorbox}
\section[\texorpdfstring{\gu{ઇન્ટરવ્યૂ} (inṭaravyū)}{ઇન્ટરવ્યૂ (inṭaravyū)}]{\gu{ઇન્ટરવ્યૂ} (inṭaravyū) — a job interview}

\label{lesson:GU-C512-intaravyu}



\begin{quote}
Before the new one: say the Gujarati for a clerk, then the Gujarati for a job post.
\end{quote}

\subsection*{You'll want to know: \gu{ઇન્ટરવ્યૂ}}
\textbf{\gu{ઇન્ટરવ્યૂ}} — \emph{inṭaravyū} — \textquotedblleft{}a job interview\textquotedblright{}.

\begin{quote}
\gu{મારો} \gu{ઇન્ટરવ્યૂ} \gu{કાલે} \gu{સવારે} \gu{છે}.
\end{quote}

\begin{quote}
\emph{māro inṭaravyū kāle savāre che.}
\end{quote}

\begin{quote}
My interview is tomorrow morning.
\end{quote}

A masculine noun.

\begin{cousinweb}
\textbf{\gu{ઇન્ટરવ્યૂ}} (\emph{inṭaravyū}) is borrowed from English \emph{interview}.
\end{cousinweb}

\begin{grammarlens}[title={What you've built}]
One more word for describing what you lived and felt.
\end{grammarlens}

\subsection*{Your turn}
\begin{itemize}
\raggedright
  \item \emph{Say it:} \emph{inṭaravyū}
  \item \emph{Say it:} \emph{inṭaravyū}, once more
  \item \emph{Say it:} \emph{māro inṭaravyū kāle savāre che.} — My interview is tomorrow morning.
  \item \emph{Recall:} say the Gujarati for a clerk, then the Gujarati for a job post, then say \emph{inṭaravyū} again
\end{itemize}

\begin{tcolorbox}[breakable,colback=teal!4,colframe=teal!35!black,title={Before you move on}]
What does \gu{ઇન્ટરવ્યૂ} mean? (\textquotedblleft{}A job interview\textquotedblright{}.) Say it once more.
\end{tcolorbox}
\section[\texorpdfstring{\gu{વેતન} (vetan)}{વેતન (vetan)}]{\gu{વેતન} (vetan) — a wage}

\label{lesson:GU-C512-vetan}



\begin{quote}
Before the new one: say the Gujarati for a job post, then the Gujarati for a job interview.
\end{quote}

\subsection*{You'll want to know: \gu{વેતન}}
\textbf{\gu{વેતન}} — \emph{vetan} — \textquotedblleft{}a wage\textquotedblright{}.

\begin{quote}
\gu{આ} \gu{કારખાનામાં} \gu{મજૂરોનું} \gu{વેતન} \gu{ઓછું} \gu{છે}.
\end{quote}

\begin{quote}
\emph{ā kārkhānāmā̃ majūronũ vetan ochũ che.}
\end{quote}

\begin{quote}
The workers' wages in this factory are low.
\end{quote}

A neuter noun.

\begin{cousinweb}
\textbf{\gu{વેતન}} (\emph{vetan}) is a tatsama, taken from Sanskrit \emph{vetana}, \textquotedblleft{}wages, hire, salary\textquotedblright{} (Monier-Williams). It is the formal word beside the everyday \textbf{\gu{પગાર}}.
\end{cousinweb}

\begin{grammarlens}[title={What you've built}]
One more word for describing what you lived and felt.
\end{grammarlens}

\subsection*{Your turn}
\begin{itemize}
\raggedright
  \item \emph{Say it:} \emph{vetan}
  \item \emph{Say it:} \emph{vetan}, once more
  \item \emph{Say it:} \emph{ā kārkhānāmā̃ majūronũ vetan ochũ che.} — The workers' wages in this factory are low.
  \item \emph{Recall:} say the Gujarati for a job post, then the Gujarati for a job interview, then say \emph{vetan} again
\end{itemize}

\begin{tcolorbox}[breakable,colback=teal!4,colframe=teal!35!black,title={Before you move on}]
What does \gu{વેતન} mean? (\textquotedblleft{}A wage\textquotedblright{}.) Say it once more.
\end{tcolorbox}
\section[\texorpdfstring{\gu{કારકિર્દી} (kārkirdī)}{કારકિર્દી (kārkirdī)}]{\gu{કારકિર્દી} (kārkirdī) — a career}

\label{lesson:GU-C512-karkirdi}



\begin{quote}
Before the new one: say the Gujarati for a job interview, then the Gujarati for a wage.
\end{quote}

\subsection*{You'll want to know: \gu{કારકિર્દી}}
\textbf{\gu{કારકિર્દી}} — \emph{kārkirdī} — \textquotedblleft{}a career\textquotedblright{}.

\begin{quote}
\gu{મીરાએ} \gu{સંગીતમાં} \gu{કારકિર્દી} \gu{બનાવી}.
\end{quote}

\begin{quote}
\emph{Mīrāe sangītmā̃ kārkirdī banāvī.}
\end{quote}

\begin{quote}
Mira made a career in music.
\end{quote}

A feminine noun.

\begin{cousinweb}
\textbf{\gu{કારકિર્દી}} (\emph{kārkirdī}) is usually traced to Persian \emph{kārkird}, \textquotedblleft{}work done, a record of deeds\textquotedblright{}, built on \emph{kār}, \textquotedblleft{}work\textquotedblright{}.
\end{cousinweb}

\begin{grammarlens}[title={What you've built}]
One more word for describing what you lived and felt.
\end{grammarlens}

\subsection*{Your turn}
\begin{itemize}
\raggedright
  \item \emph{Say it:} \emph{kārkirdī}
  \item \emph{Say it:} \emph{kārkirdī}, once more
  \item \emph{Say it:} \emph{Mīrāe sangītmā̃ kārkirdī banāvī.} — Mira made a career in music.
  \item \emph{Recall:} say the Gujarati for a job interview, then the Gujarati for a wage, then say \emph{kārkirdī} again
\end{itemize}

\begin{tcolorbox}[breakable,colback=teal!4,colframe=teal!35!black,title={Before you move on}]
What does \gu{કારકિર્દી} mean? (\textquotedblleft{}A career\textquotedblright{}.) Say it once more.
\end{tcolorbox}
\section[\texorpdfstring{\gu{બઢતી} (baḍhtī)}{બઢતી (baḍhtī)}]{\gu{બઢતી} (baḍhtī) — a promotion (at work)}

\label{lesson:GU-C512-badhti}



\begin{quote}
Before the new one: say the Gujarati for a wage, then the Gujarati for a career.
\end{quote}

\subsection*{You'll want to know: \gu{બઢતી}}
\textbf{\gu{બઢતી}} — \emph{baḍhtī} — \textquotedblleft{}a promotion (at work)\textquotedblright{}.

\begin{quote}
\gu{અરુણને} \gu{આ} \gu{વરસે} \gu{બઢતી} \gu{મળી}.
\end{quote}

\begin{quote}
\emph{Aruṇne ā varse baḍhtī maḷī.}
\end{quote}

\begin{quote}
Arun got a promotion this year.
\end{quote}

A feminine noun.

\begin{cousinweb}
\textbf{\gu{બઢતી}} (\emph{baḍhtī}) is probably borrowed from Hindi \emph{baṛhtī}, \textquotedblleft{}an increase\textquotedblright{}, built on \emph{baṛhnā}, \textquotedblleft{}to increase\textquotedblright{}, which CDIAL lists under Sanskrit \emph{várdhatē}, \textquotedblleft{}grows\textquotedblright{} (CDIAL 11376). CDIAL gives Gujarati's own verb from the same root, \textbf{\gu{વધવું}}, there too.
\end{cousinweb}

\begin{grammarlens}[title={What you've built}]
One more word for describing what you lived and felt.
\end{grammarlens}

\subsection*{Your turn}
\begin{itemize}
\raggedright
  \item \emph{Say it:} \emph{baḍhtī}
  \item \emph{Say it:} \emph{baḍhtī}, once more
  \item \emph{Say it:} \emph{Aruṇne ā varse baḍhtī maḷī.} — Arun got a promotion this year.
  \item \emph{Recall:} say the Gujarati for a wage, then the Gujarati for a career, then say \emph{baḍhtī} again
\end{itemize}

\begin{tcolorbox}[breakable,colback=teal!4,colframe=teal!35!black,title={Before you move on}]
What does \gu{બઢતી} mean? (\textquotedblleft{}A promotion (at work)\textquotedblright{}.) Say it once more.
\end{tcolorbox}
